\begin{table}[t]
\tabsize
\caption{Transfer between the three alternatives with a self-referenced GBDT detector, three-cycle windows, common features}\label{tab:transfer3}
\begin{tabular}{@{}llcccc@{}}
\toprule
Trained on & Tested on & AUC all & AUC o.o.r. & TPR o.o.r. & FA healthy (\%) \\
\midrule
PMSG & PMSG & 0.86 & 0.85 & 0.62 & 3 \\
SCIG & SCIG & 0.92 & 0.91 & 0.76 & 1 \\
WFSG & WFSG & 0.97 & 0.97 & 0.93 & -- \\
PMSG & SCIG & 0.88 & 0.88 & 0.60 & 0 \\
PMSG & WFSG & 0.97 & 0.98 & 0.89 & -- \\
SCIG & PMSG & 0.81 & 0.80 & 0.35 & 0 \\
SCIG & WFSG & 0.95 & 0.96 & 0.86 & -- \\
WFSG & PMSG & 0.86 & 0.84 & 0.52 & 1 \\
WFSG & SCIG & 0.91 & 0.89 & 0.61 & 0 \\
SCIG+WFSG & PMSG & 0.84 & 0.83 & 0.41 & 0 \\
PMSG+WFSG & SCIG & 0.89 & 0.88 & 0.61 & 0 \\
PMSG+SCIG & WFSG & 0.97 & 0.97 & 0.88 & -- \\
\botrule
\end{tabular}
\footnotetext{Features $|z|$-referenced to the first half of each trial's own reference interval. AUC on all windows and on the out-of-reference (o.o.r.) windows; TPR: true-positive rate at 1\,\% false-positive rate (FPR) on the o.o.r. windows; FA: share of the healthy trials' fault-time windows above the threshold giving 1\,\% FPR on all negatives. The WFSG has no healthy trials, so its negatives are reference and recovery windows only and its rows are optimistic.}
\end{table}
